\documentclass[11pt]{article}
\usepackage[margin=1in]{geometry}
\usepackage{amsmath,amssymb}
\usepackage{hyperref}

\newcommand{\R}{\mathbb R}

\title{Literature status for arXiv:1805.10203: Gaussian multi-bubble conjecture}
\author{Run \texttt{fa\_banach\_001}, agent \texttt{agent\_lane\_00}}
\date{2026-07-18}

\begin{document}
\maketitle

\section*{Status}

This is a lightweight literature-status packet.  The extracted target from
Steven Heilman's \emph{The Structure of Gaussian Minimal Bubbles}
\cite{Heilman2018} is already answered by separate work of Milman and Neeman
\cite{MilmanNeeman2018,MilmanNeeman2022}.  It is not a new proof or partial
result from this run.

\section*{Source Question}

The source paper states the Gaussian Multi-Bubble Problem as Problem 1.2 and
the Gaussian Multi-Bubble Conjecture as Conjecture 1.3.  In source notation,
for \(m\ge3\) and prescribed positive Gaussian volumes
\(a_1,\ldots,a_m\) summing to \(1\), Problem 1.2 asks for a measurable
partition
\[
    \Omega_1,\ldots,\Omega_m\subset\R^{n+1},
    \qquad \gamma_{n+1}(\Omega_i)=a_i,
\]
minimizing total Gaussian interface area
\[
    \sum_{1\le i<j\le m}
    \int_{\partial\Omega_i\cap\partial\Omega_j}\gamma_n(x)\,dx .
\]
Conjecture 1.3 predicts that every minimizer is, up to translation, the
Voronoi partition determined by the vertices of a centered regular simplex:
\[
    \Omega_i
    =
    y+\{x\in\R^{n+1}:\langle x,z_i\rangle
       =\max_{1\le j\le m}\langle x,z_j\rangle\}.
\]

\section*{Supporting Answer}

Milman and Neeman's supporting paper proves the Gaussian Multi-Bubble
Theorem.  In the arXiv version, Theorem 1.1 states that the Gaussian
Multi-Bubble Conjecture holds for all \(2\le q\le n+1\), and Theorem 1.2
states uniqueness up to null sets of the corresponding simplicial
\(q\)-clusters for prescribed positive Gaussian measures.

This directly answers the source's Conjecture 1.3 in the standard dimension
range where \(m=q\le n+1\) regular simplex clusters exist.  The source also
explicitly records this provenance: after Heilman's first arXiv upload,
Milman and Neeman uploaded arXiv:1805.10961, solving Conjecture 1.3
unconditionally.

\section*{Search And Evidence}

Cheap run-index search for arXiv:1805.10203, Gaussian minimal bubbles,
Gaussian Multi-Bubble Problem, Gaussian Multi-Bubble Conjecture, simplicial
cluster, and plurality-is-stablest found no existing packet.  Web/literature
search found the Annals of Mathematics publication page for Milman--Neeman,
which states that the paper establishes the Gaussian Multi-Bubble Conjecture
and gives the DOI \href{https://doi.org/10.4007/annals.2022.195.1.2}
{10.4007/annals.2022.195.1.2}.  The local source bibliography cites
arXiv:1805.10961 as \emph{The Gaussian multi-bubble conjecture}.

\section*{Scope Limitation}

This packet only records the literature status of the Gaussian multi-bubble
problem extracted from arXiv:1805.10203.  It does not address Euclidean,
spherical, hyperbolic, anisotropic, or stability variants of multi-bubble
isoperimetry unless those variants are explicitly covered by the cited
Milman--Neeman theorem.

\begin{thebibliography}{9}
\bibitem{Heilman2018}
Steven Heilman,
\emph{The Structure of Gaussian Minimal Bubbles},
arXiv:1805.10203, 2018.

\bibitem{MilmanNeeman2018}
Emanuel Milman and Joe Neeman,
\emph{The Gaussian Multi-Bubble Conjecture},
arXiv:1805.10961, 2018.

\bibitem{MilmanNeeman2022}
Emanuel Milman and Joe Neeman,
\emph{The Gaussian Double-Bubble and Multi-Bubble Conjectures},
Annals of Mathematics 195 (2022), no. 1, 89--206.
DOI: \href{https://doi.org/10.4007/annals.2022.195.1.2}
{10.4007/annals.2022.195.1.2}.
\end{thebibliography}

\end{document}
